\chapter{Distributions and Distributional Derivatives}
\label{ch:iii16}

Distributions extend the derivative and transform calculus to nonsmooth objects by duality. A distribution acts continuously and linearly on test functions, and derivatives are transferred onto the test function by integration by parts.

\section*{Learning goals}
After this chapter, the reader should be able to:
\begin{itemize}
\item verify that a linear functional defines a distribution by checking continuity on test functions;
\item embed locally integrable functions into distributions and compute their action on test functions;
\item compute distributional derivatives by transferring derivatives to test functions with the correct sign;
\item derive the derivative of the Heaviside function and related jump distributions;
\item prove convergence of distributions by testing against an arbitrary fixed test function;
\item apply smooth multiplication and differential operators to distributions without reverting to pointwise intuition;
\end{itemize}
\section*{Conceptual roadmap}
\[
\boxed{\text{structure}\;\longrightarrow\;\text{estimate}\;\longrightarrow\;\text{limit or representation}\;\longrightarrow\;\text{application}.}
\]

\paragraph{Mathematical setting.}
All distributions are continuous linear (not conjugate-linear) functionals on $\mathcal D(\Omega)$, where $\Omega\subset\mathbb R^n$ is open; $T_f(\phi)=\int_\Omega f\phi$ for $f\in L^1_{\rm loc}$. Point masses require $x_0\in\Omega$. The one-dimensional examples are on $\mathbb R$. A mollifier means $\rho\in C_c^\infty$, $\rho\ge0$, $\int\rho=1$. In jump formulas, assume finitely many jumps on each compact interval, one-sided C1 extensions at each jump, and locally integrable piecewise derivative; the sum of jump masses is then locally finite.

\section{Distributions}
A distribution on $\Omega$ is a continuous linear functional on $C_c^\infty(\Omega)$.

\section{Regular distributions}
Locally integrable functions define distributions by integration against test functions.

\section{Dirac mass}
Point evaluation $\delta_{x_0}(\phi)=\phi(x_0)$ is a distribution not represented by an ordinary locally integrable function.


% BEGIN VOL03-EXPANSION III16-example-01
\begin{example}[Dirac is continuous on test functions]\label{ex:iii16-expand-01}
For a fixed \(x_0\), define
\[
\delta_{x_0}(\phi)=\phi(x_0).
\]
If test functions \(\phi_j\to0\) with supports in one compact set, then
\[
|\delta_{x_0}(\phi_j)|\le \|\phi_j\|_\infty\to0.
\]
Thus point evaluation is a continuous linear functional and hence a
distribution. It is singular because no locally integrable function can
represent exact point evaluation against every test function.
\end{example}
% END VOL03-EXPANSION III16-example-01
\section{Distributional derivatives}
Define $\partial^\alpha T(\phi)=(-1)^{|\alpha|}T(\partial^\alpha\phi)$.


% BEGIN VOL03-EXPANSION III16-example-02
\begin{example}[The jump of Heaviside becomes a Dirac mass]\label{ex:iii16-expand-02}
Let \(H(x)=0\) for \(x<0\) and \(H(x)=1\) for \(x>0\). For a test function
\(\phi\),
\[
H'(\phi)=-\int_0^\infty \phi'(x)\,dx=\phi(0).
\]
Therefore
\[
H'=\delta_0
\]
in the distributional sense. The derivative records the jump that the
classical derivative misses away from the single discontinuity.
\end{example}
% END VOL03-EXPANSION III16-example-02
\section{Integration by parts}
For smooth functions, distributional differentiation agrees with classical differentiation because test functions have compact support.

\section{Convergence of distributions}
A sequence converges distributionally when it converges on every test function.


% BEGIN VOL03-EXPANSION III16-example-03
\begin{example}[Mollifiers converge to Dirac distributionally]\label{ex:iii16-expand-03}
Let \(\rho_\varepsilon(x)=\varepsilon^{-n}\rho(x/\varepsilon)\) with
\(\int\rho=1\). For every test function \(\phi\),
\[
\langle \rho_\varepsilon,\phi\rangle
=\int \rho(y)\phi(\varepsilon y)\,dy
\longrightarrow \phi(0)
\]
by dominated convergence. Hence
\[
\rho_\varepsilon\to\delta_0
\]
as distributions. A singular source can therefore arise as the limit of
ordinary smooth densities.
\end{example}
% END VOL03-EXPANSION III16-example-03
\section{Multiplication by smooth functions}
If $a\in C^\infty$, define $(aT)(\phi)=T(a\phi)$.

\section{Linear differential operators}
Smooth-coefficient differential operators act naturally on distributions by combining multiplication and distributional differentiation.

\section{Core structural results}

\begin{theorem}[Agreement with classical derivatives]\label{thm:iii16-01}
If $f\in C^1$ defines a regular distribution, then its distributional derivative is the regular distribution defined by $f'$.
\begin{proof}
For every test function, $D_f'(\phi)=-\int f\phi'=\int f'\phi$ by integration by parts and vanishing boundary terms.
\end{proof}
\end{theorem}

\begin{theorem}[Derivative of Heaviside]\label{thm:iii16-02}
The distributional derivative of the Heaviside step is $\delta_0$.
\begin{proof}
For a test function $\phi$, $H'(\phi)=-\int_0^\infty\phi'(x)dx=\phi(0)$.
\end{proof}
\end{theorem}

\begin{theorem}[Continuity of differentiation]\label{thm:iii16-03}
If $T_n\to T$ in distributions, then $\partial^\alpha T_n\to\partial^\alpha T$.
\begin{proof}
Evaluate on a test function and move the derivative to that fixed test function; convergence of $T_n$ then applies directly.
\end{proof}
\end{theorem}

\section{Worked examples}

\begin{example}[Dirac mass is a continuous distribution]\label{ex:iii16-worked-01}
For \(a\in\Omega\), define \(\delta_a(\varphi)=\varphi(a)\). On any
\(\mathcal D_K\) containing \(a\),
\[
|\delta_a(\varphi)|\le \sup_K|\varphi|=p_0(\varphi),
\]
while if \(a\notin K\) the value is zero. Hence \(\delta_a\) is a distribution
of order zero.
\end{example}

\begin{example}[Regular distributions from locally integrable functions]\label{ex:iii16-worked-02}
If \(f\in L^1_{\mathrm{loc}}(\Omega)\), define
\[
T_f(\varphi)=\int_\Omega f\varphi.
\]
For \(\operatorname{supp}\varphi\subset K\),
\[
|T_f(\varphi)|\le \|f\|_{L^1(K)}\sup_K|\varphi|.
\]
Thus \(T_f\) is a distribution of order zero on each compact set.
\end{example}

\begin{example}[The derivative of Heaviside is Dirac]\label{ex:iii16-worked-03}
Let \(H(x)=\mathbf1_{(0,\infty)}(x)\). Then
\[
\langle H',\varphi\rangle
=-\int_0^\infty\varphi'(x)\,dx
=\varphi(0)
=\langle\delta_0,\varphi\rangle.
\]
Hence \(H'=\delta_0\) distributionally.
\end{example}

\begin{example}[A piecewise smooth jump creates a delta]\label{ex:iii16-worked-04}
If \(u\) is \(C^1\) away from \(0\) and has jump
\(J=u(0^+)-u(0^-)\), then its distributional derivative contains
\[
J\delta_0
\]
in addition to the classical derivative on the two sides. Integration by parts
on the two half-lines exposes the boundary term.
\end{example}

\section{Solved dossiers}
Each dossier combines several ideas from the chapter in a longer argument. The aim is to make hypotheses, calculations, and proof strategy explicit.

\begin{problem}[Dirac mass is a continuous distribution]\label{prob:iii16-dossier-01}
Prove the Dirac mass is a distribution.
\end{problem}
\begin{solution}
Evaluation is linear, and on each fixed compact support it is bounded by
the zeroth test-function seminorm.
\end{solution}

\begin{problem}[Regular distributions from locally integrable functions]\label{prob:iii16-dossier-02}
Show a locally integrable function defines a distribution.
\end{problem}
\begin{solution}
On each compact \(K\), bound the pairing by
\(\|f\|_{L^1(K)}p_0(\varphi)\).
\end{solution}

\begin{problem}[The derivative of Heaviside is Dirac]\label{prob:iii16-dossier-03}
Compute the distributional derivative of the Heaviside function.
\end{problem}
\begin{solution}
Integration by parts against a compactly supported test function gives
\(\langle H',\varphi\rangle=\varphi(0)\), so \(H'=\delta_0\).
\end{solution}

\begin{problem}[A piecewise smooth jump creates a delta]\label{prob:iii16-dossier-04}
Explain how a jump discontinuity appears in the distributional derivative.
\end{problem}
\begin{solution}
Piecewise integration by parts yields the classical derivative plus a Dirac
mass whose coefficient equals the jump size.
\end{solution}

\begin{problem}[Derivative of Dirac]\label{prob:iii16-dossier-05}
Define \(\delta_0'\) and state its action on a test function.
\end{problem}
\begin{solution}
\(\langle\delta_0',\varphi\rangle=-\varphi'(0)\).
\end{solution}

\begin{problem}[Distributional derivative definition]\label{prob:iii16-dossier-06}
State the multi-index derivative of a distribution.
\end{problem}
\begin{solution}
\(\langle D^\alpha T,\varphi\rangle=(-1)^{|\alpha|}
\langle T,D^\alpha\varphi\rangle\).
\end{solution}

\begin{problem}[Classical compatibility]\label{prob:iii16-dossier-07}
If \(f\in C^1\), show its distributional derivative agrees with its classical derivative.
\end{problem}
\begin{solution}
Integrate by parts; the boundary term vanishes because test functions are
compactly supported.
\end{solution}

\begin{problem}[Continuity is essential]\label{prob:iii16-dossier-08}
Why is a distribution not merely an arbitrary linear functional on \(C_c^\infty\)?
\end{problem}
\begin{solution}
The definition requires continuity in the test-function topology; without
it, pathological algebraic linear functionals would be included and limit
operations would fail.
\end{solution}

\begin{problem}[Agreement with classical derivatives proof dossier]\label{prob:iii16-dossier-09}
Prove the following structural statement and identify the step where its main hypothesis is used: If $f\in C^1$ defines a regular distribution, then its distributional derivative is the regular distribution defined by $f'$.
\end{problem}
\begin{solution}
For every test function, $D_f'(\phi)=-\int f\phi'=\int f'\phi$ by integration by parts and vanishing boundary terms. This proof also explains why the stated hypothesis is structural rather than cosmetic.
\end{solution}

\begin{problem}[Derivative of Heaviside proof dossier]\label{prob:iii16-dossier-10}
Prove the following structural statement and identify the step where its main hypothesis is used: The distributional derivative of the Heaviside step is $\delta_0$.
\end{problem}
\begin{solution}
For a test function $\phi$, $H'(\phi)=-\int_0^\infty\phi'(x)dx=\phi(0)$. This proof also explains why the stated hypothesis is structural rather than cosmetic.
\end{solution}

\begin{problem}[Continuity of differentiation proof dossier]\label{prob:iii16-dossier-11}
Prove the following structural statement and identify the step where its main hypothesis is used: If $T_n\to T$ in distributions, then $\partial^\alpha T_n\to\partial^\alpha T$.
\end{problem}
\begin{solution}
Evaluate on a test function and move the derivative to that fixed test function; convergence of $T_n$ then applies directly. This proof also explains why the stated hypothesis is structural rather than cosmetic.
\end{solution}

\begin{problem}[Chapter synthesis]\label{prob:iii16-dossier-12}
Connect the main definitions and structural theorems of this chapter into one proof strategy. Explain what object is controlled, what estimate or closure principle is used, and what conclusion becomes available.
\end{problem}
\begin{solution}
Distributions extend the derivative and transform calculus to nonsmooth objects by duality. A distribution acts continuously and linearly on test functions, and derivatives are transferred onto the test function by integration by parts. The recurring strategy is to encode the object in the chapter's natural measurable, normed, Fourier, or distributional structure, establish the controlling estimate, and then pass to the desired limit or representation.
\end{solution}
\section{Exercises with complete solutions}
The shorter exercise layer remains separate from the solved dossiers.

\begin{exercise}\label{exr:iii16-01}
State and apply the central principle behind Distributions. Give one immediate consequence in the setting of this chapter.
\end{exercise}
\begin{hint}
To verify a distribution, check linearity and estimate its action by finitely many test-function seminorms on each compact set.
\end{hint}
\begin{solution}
A distribution on $\Omega$ is a continuous linear functional on $C_c^\infty(\Omega)$. As an immediate consequence, the same principle can be inserted into later convergence, approximation, transform, or weak-form arguments whenever its hypotheses are verified.
\end{solution}

\begin{exercise}\label{exr:iii16-02}
State and apply the central principle behind Regular distributions. Give one immediate consequence in the setting of this chapter.
\end{exercise}
\begin{hint}
A locally integrable function acts by \(T_f(\phi)=\int f\phi\); use local integrability because \(\phi\) has compact support.
\end{hint}
\begin{solution}
Locally integrable functions define distributions by integration against test functions. As an immediate consequence, the same principle can be inserted into later convergence, approximation, transform, or weak-form arguments whenever its hypotheses are verified.
\end{solution}

\begin{exercise}\label{exr:iii16-03}
State and apply the central principle behind Dirac mass. Give one immediate consequence in the setting of this chapter.
\end{exercise}
\begin{hint}
The distributional derivative is defined by \(T'(\phi)=-T(\phi')\); the minus sign comes from integration by parts.
\end{hint}
\begin{solution}
Point evaluation $\delta_{x_0}(\phi)=\phi(x_0)$ is a distribution not represented by an ordinary locally integrable function. As an immediate consequence, the same principle can be inserted into later convergence, approximation, transform, or weak-form arguments whenever its hypotheses are verified.
\end{solution}

\begin{exercise}\label{exr:iii16-04}
State and apply the central principle behind Distributional derivatives. Give one immediate consequence in the setting of this chapter.
\end{exercise}
\begin{hint}
For Heaviside, integrate \(-\int_0^\infty \phi'(x)\,dx\) and use compact support to evaluate the endpoint.
\end{hint}
\begin{solution}
Define $\partial^\alpha T(\phi)=(-1)^{|\alpha|}T(\partial^\alpha\phi)$. As an immediate consequence, the same principle can be inserted into later convergence, approximation, transform, or weak-form arguments whenever its hypotheses are verified.
\end{solution}

\begin{exercise}\label{exr:iii16-05}
State and apply the central principle behind Integration by parts. Give one immediate consequence in the setting of this chapter.
\end{exercise}
\begin{hint}
To prove convergence of distributions, fix an arbitrary test function first and pass to the scalar limit of the pairings.
\end{hint}
\begin{solution}
For smooth functions, distributional differentiation agrees with classical differentiation because test functions have compact support. As an immediate consequence, the same principle can be inserted into later convergence, approximation, transform, or weak-form arguments whenever its hypotheses are verified.
\end{solution}

\begin{exercise}\label{exr:iii16-06}
State and apply the central principle behind Convergence of distributions. Give one immediate consequence in the setting of this chapter.
\end{exercise}
\begin{hint}
Multiplication by a smooth function is defined by moving the multiplier onto the test function: \((aT)(\phi)=T(a\phi)\).
\end{hint}
\begin{solution}
A sequence converges distributionally when it converges on every test function. As an immediate consequence, the same principle can be inserted into later convergence, approximation, transform, or weak-form arguments whenever its hypotheses are verified.
\end{solution}

\begin{exercise}\label{exr:iii16-07}
State and apply the central principle behind Multiplication by smooth functions. Give one immediate consequence in the setting of this chapter.
\end{exercise}
\begin{hint}
For a differential operator, move every derivative onto the test function with the appropriate sign before simplifying.
\end{hint}
\begin{solution}
If $a\in C^\infty$, define $(aT)(\phi)=T(a\phi)$. As an immediate consequence, the same principle can be inserted into later convergence, approximation, transform, or weak-form arguments whenever its hypotheses are verified.
\end{solution}

\begin{exercise}\label{exr:iii16-08}
State and apply the central principle behind Linear differential operators. Give one immediate consequence in the setting of this chapter.
\end{exercise}
\begin{hint}
Do not evaluate a general distribution at a point; only its action on test functions is defined without extra structure.
\end{hint}
\begin{solution}
Smooth-coefficient differential operators act naturally on distributions by combining multiplication and distributional differentiation. As an immediate consequence, the same principle can be inserted into later convergence, approximation, transform, or weak-form arguments whenever its hypotheses are verified.
\end{solution}


% BEGIN VOL03-EXPANSION III16-exercises-01
\section{Graded supplementary exercises}
These exercises extend the preserved foundational set and are graded by mathematical role.

\subsection*{Standard computations}

\begin{exercise}[Dirac action]\label{exr:iii16-expand-01}
What is delta\_a(phi)?
\end{exercise}
\begin{hint}
Use point evaluation.
\end{hint}
\begin{solution}
It is phi(a).
\end{solution}

\begin{exercise}[Dirac derivative]\label{exr:iii16-expand-02}
What is delta'\_0(phi)?
\end{exercise}
\begin{hint}
Use the distributional derivative definition.
\end{hint}
\begin{solution}
It is -phi'(0).
\end{solution}

\begin{exercise}[Heaviside derivative]\label{exr:iii16-expand-03}
What is the distributional derivative of the Heaviside step?
\end{exercise}
\begin{hint}
Integrate by parts on the positive half-line.
\end{hint}
\begin{solution}
delta\_0.
\end{solution}

\begin{exercise}[Smooth multiplier]\label{exr:iii16-expand-04}
How is aT defined for smooth a?
\end{exercise}
\begin{hint}
Move multiplication to the test function.
\end{hint}
\begin{solution}
(aT)(phi)=T(a phi).
\end{solution}

\begin{exercise}[Distributional convergence]\label{exr:iii16-expand-05}
What must be checked to prove T\_n tends to T distributionally?
\end{exercise}
\begin{hint}
Fix an arbitrary test function.
\end{hint}
\begin{solution}
Show T\_n(phi) tends to T(phi) for every test phi.
\end{solution}

\subsection*{Proofs}

\begin{exercise}[Regular derivative]\label{exr:iii16-expand-06}
Prove a C1 function has the same classical and distributional derivative.
\end{exercise}
\begin{hint}
Integrate by parts.
\end{hint}
\begin{solution}
Compact support kills the boundary term.
\end{solution}

\begin{exercise}[Derivative continuity]\label{exr:iii16-expand-07}
Prove differentiation is continuous on distributions.
\end{exercise}
\begin{hint}
Move the derivative to the fixed test function.
\end{hint}
\begin{solution}
Convergence of T\_n on that derivative gives convergence of derivatives.
\end{solution}

\begin{exercise}[Linearity]\label{exr:iii16-expand-08}
Prove the distributional derivative is linear.
\end{exercise}
\begin{hint}
Apply the definition to a linear combination.
\end{hint}
\begin{solution}
The sign and test derivative distribute linearly.
\end{solution}

\begin{exercise}[Leibniz smooth factor]\label{exr:iii16-expand-09}
Derive D(aT)=a DT+a' T in one dimension.
\end{exercise}
\begin{hint}
Evaluate on a test function and expand (a phi)'.
\end{hint}
\begin{solution}
The two terms arise from the product rule.
\end{solution}

\subsection*{Counterexamples and hypothesis testing}

\begin{exercise}[Dirac not ordinary function]\label{exr:iii16-expand-10}
Why can no locally integrable f represent delta\_0?
\end{exercise}
\begin{hint}
Test with bumps supported in shrinking neighborhoods of zero.
\end{hint}
\begin{solution}
Absolute continuity of the integral would force the actions to shrink, while delta can stay one.
\end{solution}

\begin{exercise}[Pointwise derivative failure]\label{exr:iii16-expand-11}
Why is H'=0 away from zero not the full derivative in distributions?
\end{exercise}
\begin{hint}
Distributional differentiation tests across the jump.
\end{hint}
\begin{solution}
The missing jump contributes delta\_0.
\end{solution}

\begin{exercise}[Weak convergence versus pointwise]\label{exr:iii16-expand-12}
Why does distributional convergence not imply pointwise convergence?
\end{exercise}
\begin{hint}
Distributions may be singular and have no point values.
\end{hint}
\begin{solution}
For example $f_n(x)=\sin(nx)$ converges distributionally to zero on $\mathbb R$ by integration by parts against each compactly supported smooth test. At $x=\pi/2$ its values cycle through $1,0,-1,0$, so pointwise convergence fails.
\end{solution}

\subsection*{Applications and investigations}

\begin{exercise}[Impulse model]\label{exr:iii16-expand-13}
Why is Dirac useful for point forces or instantaneous impulses?
\end{exercise}
\begin{hint}
It concentrates finite action at a lower-dimensional set.
\end{hint}
\begin{solution}
Differential equations can encode idealized sources while remaining linear in the distribution framework.
\end{solution}

\begin{exercise}[Approximation of sources]\label{exr:iii16-expand-14}
Why replace delta by a narrow mollifier numerically?
\end{exercise}
\begin{hint}
It provides a smooth regularization with the same total mass.
\end{hint}
\begin{solution}
The regularized solutions can converge distributionally to the singular-source solution.
\end{solution}

\subsection*{Challenge problems}

\begin{exercise}[Jump formula]\label{exr:iii16-expand-15}
If a piecewise C1 function has jumps J\_k at points x\_k, describe its distributional derivative.
\end{exercise}
\begin{hint}
Differentiate classically on smooth pieces and collect internal boundary terms.
\end{hint}
\begin{solution}
It is the classical piecewise derivative plus sum J\_k delta\_{x\_k}.
\end{solution}

\begin{exercise}[Second derivative of |x|]\label{exr:iii16-expand-16}
Compute the second distributional derivative of |x|.
\end{exercise}
\begin{hint}
First derivative is sign(x); its jump at zero has size 2.
\end{hint}
\begin{solution}
The second derivative is 2 delta\_0.
\end{solution}
% END VOL03-EXPANSION III16-exercises-01
\section*{Chapter summary}
The chapter now supplies a canonical theorem-and-dossier layer for subsequent measure, Fourier, distributional, Sobolev, and PDE arguments.
